\documentclass[10pt]{article} 

\usepackage[a4paper,margin=2cm,left=2cm,right=2cm]{geometry} 
\usepackage[UKenglish]{babel}

\usepackage{amsmath}
\usepackage{amssymb}
\usepackage{amsthm}

\newtheorem{theorem}{Theorem}
\newtheorem{lemma}{Lemma}
\newtheorem{corollary}{Corollary}[theorem]
\newtheorem{proposition}{Proposition}
\theoremstyle{definition}
\newtheorem{definition}{Definition}
\newtheorem{example}{Example}
\theoremstyle{remark}
\newtheorem{remark}{Remark}

\newcommand{\R}{\mathbb{R}}
\newcommand{\Z}{\mathbb{Z}}
\newcommand{\N}{\mathbb{N}}
\newcommand{\Q}{\mathbb{Q}}

% Example
% \DeclareMathOperator{\dom}{dom}

\usepackage{graphicx}
\graphicspath{{./images/}}

% UNCOMMENT IF USING REFERENCES
% \usepackage[notes, backend=biber]{biblatex-chicago}
% \addbibresource{refs.bib} 

\title{Lean Reference}
\author{joel}
\date{}

\begin{document}

\maketitle

\section{Tactics}

Almost everything that you will do in Lean is done through tactics if you are using 
it as a proof assistant as opposed to a pure functional programming language. 

\verb|rfl| (reflexivity (wrt equality)) is the tactic that states that two
things are equal. 

\verb|rw [h]| (rewrite) is the tactic that given a theorem \verb|h: X = Y| we
make the substituion of all $X$ to $Y$. Adding an arrow before the $h$,
(denoted by $\l$) gives the opposite, i.e turns all $Y$ into $X$. 

Can also do \verb|repeat rw [h]| in order to do it multiple times, by default
will only convert the first instance of it. 



% \printbibliography

\end{document}
